\documentclass[10pt,letterpaper]{article}
\usepackage[T1]{fontenc}
\usepackage{lmodern,amsmath,amssymb,amsthm,microtype}
\usepackage[margin=0.9in]{geometry}
\usepackage[hidelinks]{hyperref}
\newtheorem{theorem}{Theorem}
\title{A short-gap normal form for separated subset sums}
\author{Aafi Khan Malek\\ \small Independent researcher, Ahmedabad, India\\
\small\href{mailto:aafimalek2032@gmail.com}{aafimalek2032@gmail.com}}
\date{28 November 2026}
\begin{document}
\maketitle
\begin{abstract}
For sorted positive real weights with pairwise unit-separated
subset sums, all adjacent gaps below two lie in one pair or one
triple. That pair or triple can be compressed to unit gaps at
fixed mean without changing the order of any labeled subset sums
or increasing the largest weight. For nine weights below 150,
this reduces the last slice to a terminal pair or triple with
unit gaps. A separate midpoint estimate gives a new sixth-weight
lower bound. Exact rational replay then excludes two more central
sign orientations, leaving 122 unresolved.
\end{abstract}

\begin{theorem}\label{normal}
Let \(0<x_1<\cdots<x_n\) have pairwise unit-separated subset sums.
There is at most one nontrivial block of adjacent gaps below two,
and it contains at most three weights. If there is such a block,
replace its weights by consecutive unit-spaced weights with the same
arithmetic mean, and leave the others fixed. The new weights
are positive, the largest does not increase, all subset sums
remain unit-separated, and the strict order of the labeled sums
is unchanged. Every new adjacent gap is one or at least two.
\end{theorem}
\begin{proof}
All adjacent gaps are at least one. Two gaps at indices
\(i<j-1\), both below two, would give
\[
 \left|(x_{i+1}+x_j)-(x_i+x_{j+1})\right|<1,
\]
from their two lengths in \([1,2)\). The displayed sums are
two distinct labeled subsets. This contradiction leaves at most
two short gaps, necessarily consecutive if there are two. If none
exists, no replacement is needed.

To analyze a replacement, decompose all subset sums into translates
of the scores of weights outside the block. A positive difference
\(t\) of two outside scores satisfies \(t\ge1\) and
\(|t-p|\ge1\) for each positive difference \(p\) of two block
scores, by comparing the two opposite translations. We will
keep the signs of these comparisons and the unit distances,
as well as the labeled order and spacing for block scores.

For a two-weight block \(X<M\), put \(d=M-X\) and \(T=X+M\).
The positive block differences are \(d,X,M,T\).
Since \(d<2\), \(t\ge1\) and \(|t-d|\ge1\) force
\(t\ge1+d\). The forbidden intervals of radius one around
\(X,M\) overlap, so \(t\le X-1\) or \(t\ge M+1\).
Replace the two weights by \((T-1)/2,(T+1)/2\).
The new gap is one; the two new singleton differences lie
between \(X,M\); and \(T\) is fixed. These changes preserve
every required sign and distance relative to \(t\).
The four block scores remain in the same order with gaps
at least one.

For a three-weight block \(Y<X<M\), put \(e=X-Y\),
\(d=M-X\), \(T=Y+X+M\), and \(C=T/3\).
Both gaps are in \([1,2)\).
The comparison of \(Y+X\) with \(M\) gives \(|Y-d|\ge1\).
Since \(Y\ge1\) and \(d<2\), we have \(Y-d>-1\), hence
\(Y-d\ge1\). In particular \(C\ge3\).
Replace the block by \(C-1,C,C+1\). These lie in \([Y,M]\):
\(C-1\ge Y\) and \(C+1\le M\) follow by averaging the two
adjacent gaps, each at least one. The new block scores are
\[
 0,\ C-1,\ C,\ C+1,\ 2C-1,\ 2C,\ 2C+1,\ 3C,
\]
in the same label order as before, with all gaps at least one.

The positive differences of block scores are exhausted, up to
repetition, by four classes of signed support:
\begin{itemize}
\item The positive single-weight differences \(d,e,d+e\).
Their overlapping radius-one intervals cover 1 and reach
\(d+e+1\), so \(t\ge d+e+1\ge3\). Their replacements are 1 or 2.
\item The three singleton scores and the three differences of
a two-weight score and its remaining singleton. They lie
between \(Y-d\) and \(M+e\). The five centers
\(Y-d,Y,X,M,M+e\) have gaps \(d,e,d,e<2\), so the whole
range with radius-one endpoints is forbidden to \(t\).
Every corresponding new center lies in \([C-2,C+2]\),
inside \([X-d-e,X+d+e]\) since
\(|C-X|=|d-e|/3\le d+e-2\).
\item The three pair sums \(T-M,T-X,T-Y\).
Their forbidden intervals likewise overlap because their gaps
are \(d,e<2\). The new pair sums remain in that closed
range, since the new individual weights lie in \([Y,M]\)
and \(T\) is fixed.
\item The total \(T\), which is unchanged.
\end{itemize}
Even if the numerical ranges of these classes overlap, each
signed support keeps its prescribed comparison with \(t\)
and at least unit distance. When outside scores coincide,
the block-only check already gives unit separation.
Every full labeled comparison now has the same sign as before.
The block extrema moved inward, so external gaps cannot
decrease. This proves the theorem.
\end{proof}

\section*{The final nine-weight slice below 150}
The
\href{https://github.com/Aafimalek/novel-discoveries/tree/c582f2587588eeeb195048a74552fe80bdc1d0e7/nine-separated-subset-sums/verification/rank-cuts-below-150/only-floor-148}{previous rank-swap reduction}
shows that in the remaining slice with largest weight below 150,
\[
 149\le M=x_9<150,\qquad 148\le X=x_8<149,\qquad
 1\le d=M-X<41/39.
\]
By Theorem~\ref{normal}, we may preserve the entire order of
the 512 labeled sums and assume \(M=X+1\). The only close block
is a terminal pair or triple, so it can also be assumed that
\(X-x_7\ge2\) or \(X-x_7=1\). If the largest weight is
decreased during normalization, it stays at least 149 by
the existing lower bound, so the normalized point remains
in the required slice. Keeping the seven smaller weights fixed
and decreasing normalized \(X\) within any given feasible
chamber, its least possible value is \(148\), \(t+1\), or
\(t/2\) for a positive difference \(t\) of the first-seven scores:
the variable forbidden intervals are \((X-1,X+2)\)
and \((2X,2X+2)\), whose two branches give the listed lower
endpoints or upper bounds. This observation does not assert
that every endpoint is feasible.

\paragraph{A sixth-weight cut before normalization.}
The first seven weights are \((1+d)\)-separated.
Set \(Z=(X+M)/2\), and adjoin \(Z\) to those weights.
The new eight weights have separation \(h=1+d/2\),
since any positive first-seven difference \(t\) lies at most
\(X-1\) or at least \(M+1\), and hence \(|t-Z|\ge h\).
Replace \(x_7,Z\) by \(Z-x_7\), leaving the first six:
translating by \(x_7\) makes all resulting sums distinct
members of the eight-weight family, preserving separation.
The
\href{https://arxiv.org/abs/2611.17360}{seven-weight optimum reported by Notik}
gives a weight at least \(44h\) in the resulting seven-weight set.
But \(x_7\ge44(1+d)\ge88\), so \(Z-x_7<61.5<44h\).
Therefore \(x_6\ge44+22d\ge66\).

\paragraph{Exact screen.}
The exact rational audit covers all 124 orientations left by the
\href{https://github.com/Aafimalek/novel-discoveries/tree/4091ab711e3e0e42f8a78d9c52eb19e561e69c0a/nine-separated-subset-sums/verification/rank-cuts-below-150/floor148-exact-root-screen}{previous screen}.
It adds the unit-gap equality and midpoint bounds, and proves
two new exclusions; both certificates use only \(x_6\ge66\)
among the added rows and so apply before normalization.
There are now 390 exclusions out of 512 and 122 unresolved
orientations. Neither a complete slot search nor the terminal
pair/triple split was used here. The integer search and formal
proof of 149 are separate. This written normalization is
not Lean-checked and does not prove a lower bound of 150.
\medskip\noindent
This research note was developed with AI assistance.
\end{document}
